Shape optimization of corrugated media becomes physically interpretable when material, initial envelope and geometric admissibility are compared consistently. We reformulate the design resource as actual medium volume per initial projected area and evaluate NURBS curvature at stationary points and both sides of its knots. An audit of 572 archived finite-element geometries finds 16 radius-screen violations previously admitted by smoothed curvature. Correcting only the material coordinate reduces an archived 31-member material--stress front to one member under its historical eligibility conditions. A new experiment performs 449 direct compression paths for sinusoidal, Fourier and NURBS shapes at fixed pitch, height and medium resource. Three paired 64-call comparisons show no consistent NSGA-II advantage over Sobol sampling. Ten finalists receive 270 additional paths across 27 geometric and material sensitivity scenarios. Four remain qualified; joint perturbations reject one candidate that passed all 11 axial scenarios. At 24 elements, the qualified NURBS energy extreme retains 7.91\% greater minimum stored potential per footprint than the sinusoid, with 3.38\% greater maximum stress utilization in the sampled set. The nominal minimum-stress design and normalized knee fail geometric qualification. These results establish a practical selection principle: resource matching, knot-sided curvature and scenario enrichment materially change which nominal profile trade-offs remain admissible. The findings are conditional numerical design evidence for an attached medium between rigid surfaces; stored potential is not irreversible crush dissipation.
